\section{OP--DT theorem integration patch dossier}
\label{sec:op-dt-theorem-integration-patch}

\subsection{Claim attacked}

The scalar-square block currently separates the Oberdieck--Pixton
theorem from the Hilbert-scheme Donaldson--Thomas series:
\[
Z^X_{\mathrm{OP}}=-4096\Delta_5^{-2},
\qquad
Z^X_{\mathrm{DT}}=C_{\mathrm{DT}}\Delta_5^{-2}
\quad\hbox{under an assumption.}
\]
The assumption is no longer the strongest true statement.  In the
quotient-by-\(E\), Behrend-weighted normalization used in the manuscript,
Oberdieck's reduced DT/PT theorem identifies the Hilbert-scheme DT series
with the reduced stable-pairs series, and the primitive
Oberdieck--Pandharipande comparison identifies that stable-pairs series
with the Oberdieck--Pixton partition function.  Therefore
\begin{equation}
\label{eq:cell14-main}
Z^X_{\mathrm{DT}}(T,Q,P)
=Z^X_{\mathrm{OP}}(P,Q,T)
=-4096\Delta_5(Z)^{-2},
\qquad
C_{\mathrm{DT}}=-4096.
\end{equation}
This is an enumerative theorem.  It does not construct the microscopic
BPS Hilbert space \(\mathcal H_{\gamma(h,d,n)}\), nor a physical operator
product on it.

\subsection{Primary-source anchors}

\begin{center}
\small
\begin{tabular}{p{0.25\textwidth}p{0.66\textwidth}}
\hline
Source & Load-bearing assertion \\
\hline
Oberdieck, \emph{On reduced stable pair invariants},
arXiv:1605.04631, local source
\texttt{/tmp/igusa-cell8-sources/obred/reducedK3xE.tex:289--322}
&
Defines quotient stable-pairs invariants by the Behrend-weighted
orbifold Euler characteristic of \(P_n(X,(\beta,d))/E\), and proves
their equality with the incidence/reduced-virtual-class invariants. \\
Oberdieck, official bibliographic pages
\texttt{https://arxiv.org/abs/1605.04631} and
\texttt{https://link.springer.com/article/10.1007/s00209-017-1953-5}
&
The arXiv page records the 2016 source and abstract.  Springer records
the final publication as Mathematische Zeitschrift 289, pages 323--353,
published in 2018. \\
Oberdieck, arXiv:1605.04631, local source
\texttt{/tmp/igusa-cell8-sources/obred/reducedK3xE.tex:432--465}
&
Defines reduced Hilbert-scheme DT invariants
\(\mathsf{DT}_{n,(\beta,d)}
=\int_{\operatorname{Hilb}_X((\beta,d),n)/E}\nu\,de\),
defines the fixed-class series
\(\mathsf{DT}_{(\beta,d)}(q)\) and \(\mathsf{PT}_{(\beta,d)}(q)\), and
proves Theorem~2(i):
\[
\mathsf{DT}_{(\beta,d)}(q)=\mathsf{PT}_{(\beta,d)}(q).
\]
The same passage records that the degree-zero factor is absent because
\(\chi(S\times E)=0\). \\
Oberdieck--Pandharipande, arXiv:1411.1514, local source
\texttt{/tmp/igusa-cell8-sources/opand/K3TJul15.tex:1488--1523}
&
For primitive \(\beta_h\), the reduced stable-pairs series equals the
reduced Gromov--Witten series after
\[
y=-\exp(iu).
\]
The proof explicitly obtains the reduced correspondence for
\(S\times E\). \\
Oberdieck--Pixton, arXiv:1706.10100, local source
\texttt{/tmp/igusa-cell8-sources/op/main.tex:326--360,376--386}
&
Defines
\[
\mathcal Z_{\mathrm{OP}}(u,q,\widetilde q)
=\sum N_{g,h,d}u^{2g-2}q^{h-1}\widetilde q^{d-1}
\]
and proves
\[
\mathcal Z_{\mathrm{OP}}(u,q,\widetilde q)
=-\chi_{10}(p,q,\widetilde q)^{-1},
\qquad p=e^{iu}.
\]
\\
Bryan, arXiv:1504.02920, local source
\texttt{/tmp/igusa-cell8-sources/bryan/K3xE.tex:282--363}
&
Uses the same quotient-by-\(E\), Behrend-weighted Hilbert-scheme DT
normalization and the sign
\[
\widetilde q^{\,h-1}q^{d-1}(-p)^n.
\]
It records that Bryan's \(q,\widetilde q\) are opposite to
Oberdieck--Pandharipande's variables. \\
\hline
\end{tabular}
\end{center}

\subsection{Variable map}

Use the Oberdieck--Pixton variables
\[
P=e^{iu}=e^{2\pi iz},\qquad
Q=e^{2\pi i\tau},\qquad
T=e^{2\pi i\widetilde\tau}.
\]
The manuscript's Borcherds variables are
\begin{equation}
\label{eq:cell14-borch-vars}
q_{\mathrm{Bch}}=Q,\qquad r_{\mathrm{Bch}}=P,\qquad
s_{\mathrm{Bch}}=T.
\end{equation}
The displayed DT variables in \texttt{proj.tex:567--570} are
\[
Z^X_{\mathrm{DT}}(q,t,p)
=\sum_{h,d,n}
\mathbf{DT}^X_{n,(\beta_h,d)}
q^{d-1}t^{h-1}(-p)^n,
\qquad h-1=\frac12\langle\beta_h,\beta_h\rangle.
\]
Bryan's DT series uses \(\widetilde q^{\,h-1}q^{d-1}(-p)^n\) and
records that his \(q,\widetilde q\) are opposite to the
Oberdieck--Pandharipande variables.  Therefore
\begin{equation}
\label{eq:cell14-dt-vars}
q_{\mathrm{DT}}=T,\qquad
t_{\mathrm{DT}}=Q,\qquad
p_{\mathrm{DT}}=P.
\end{equation}
The stable-pairs sign is the same sign:
\[
y=-\exp(iu)=-P,
\qquad
(-p_{\mathrm{DT}})^n=(-P)^n.
\]

\subsection{Constants}

The manuscript normalization in \texttt{proj.tex:572--635} gives
\[
D_5=64^{-1}\Delta_5,
\qquad
\chi_{10}^{\mathrm{OP}}=D_5^2=4096^{-1}\Delta_5^2.
\]
Oberdieck--Pixton gives
\[
Z^X_{\mathrm{OP}}
=-\left(\chi_{10}^{\mathrm{OP}}\right)^{-1}.
\]
Consequently
\begin{equation}
\label{eq:cell14-constant}
Z^X_{\mathrm{OP}}
=-4096\Delta_5^{-2}.
\end{equation}
Oberdieck's reduced DT/PT theorem has no residual degree-zero factor for
\(X=S\times E\):
\[
\mathsf{DT}_0(x)=M(-x)^{\chi(X)}
=M(-x)^{\chi(S)\chi(E)}
=M(-x)^0
=1.
\]
Combining the DT/PT theorem, the primitive GW/PT comparison, and
\eqref{eq:cell14-constant} gives the fixed constant
\[
C_{\mathrm{DT}}=C_{\mathrm{OP}}=-4096.
\]

\subsection{Bibliography insertion}

Insert the following entry in \texttt{proj.bib} after
\texttt{proj.bib:190--196} and before the Thomas DT entry.  The entry is
the final-published bibliographic form of the 2016 arXiv source used in
the proof.  Springer records Mathematische Zeitschrift 289, pages
323--353, DOI \texttt{10.1007/s00209-017-1953-5}.

\begin{verbatim}
@ARTICLE{OberdieckReducedPT,
AUTHOR="G. Oberdieck",
TITLE="On reduced stable pair invariants",
JOURNAL="Mathematische Zeitschrift",
VOLUME="289",
NUMBER="1--2",
PAGES="323--353",
YEAR="2018",
DOI="10.1007/s00209-017-1953-5",
EPRINT="arXiv:1605.04631",
}
\end{verbatim}

\subsection{Normalization-table replacement}

Replace the scalar-square row in \texttt{proj.tex:534--537} by:
\begin{center}
\small
\begin{tabular}{p{0.25\textwidth}p{0.62\textwidth}}
\hbox{Scalar square} &
\(Z^X_{\mathrm{DT}}(T,Q,P)=Z^X_{\mathrm{OP}}(P,Q,T)
=-4096\Delta_5^{-2}\);
for the OP/DT branch,
\((\mathcal Z_{\mathrm{BPS}}^{\mathrm{vir}})^2
=-Z^X_{\mathrm{DT}}/4096\).
If an external scalar-square branch is later introduced, write it
separately as \(Z^X_{\square}=C_{\square}\Delta_5^{-2}\). \\
\end{tabular}
\end{center}

\subsection{Theorem replacement}

Replace \texttt{proj.tex:638--654} by the following theorem.  The label
\texttt{thm:op-square} is retained because downstream references already
point to it.

\begin{theorem}[Oberdieck--Pixton/Oberdieck reduced DT square]
\label{thm:op-square}
Let \(X=S\times E\), and let
\(\beta_h\in H_2(S,\mathbb Z)\) be primitive and non-zero with
\[
\langle\beta_h,\beta_h\rangle=2h-2.
\]
Define
\[
\mathbf{DT}^X_{n,(\beta_h,d)}
=\int_{\operatorname{Hilb}^n(X,(\beta_h,d))/E}\nu\,de
\]
and
\[
Z^X_{\mathrm{DT}}(q,t,p)
=\sum_{h,d,n}\mathbf{DT}^X_{n,(\beta_h,d)}
q^{d-1}t^{h-1}(-p)^n.
\]
Under the variable map
\[
p=P=e^{iu},\qquad t=Q,\qquad q=T,
\qquad
(q_{\mathrm{Bch}},r_{\mathrm{Bch}},s_{\mathrm{Bch}})=(Q,P,T),
\]
one has
\begin{equation}
\label{eq:dt-op-square}
Z^X_{\mathrm{DT}}(T,Q,P)
=Z^X_{\mathrm{OP}}(P,Q,T)
=-4096\,\Delta_5(Z)^{-2}.
\end{equation}
Equivalently,
\[
C_{\mathrm{DT}}=C_{\mathrm{OP}}=-4096.
\]
\end{theorem}

\begin{proof}
Oberdieck's reduced DT/PT theorem for \(S\times E\),
\cite[Theorem 3(i)]{OberdieckReducedPT}, gives, for every non-zero
\(\beta\in H_2(S,\mathbb Z)\),
\[
\sum_n\mathbf{DT}^X_{n,(\beta,d)}x^n
=
\sum_n\mathsf P^X_{n,(\beta,d)}x^n.
\]
Set \(x=-P\).  The degree-zero factor is
\[
M(-P)^{\chi(S\times E)}=M(-P)^0=1.
\]
For primitive \(\beta_h\), the reduced GW/pairs comparison of
Oberdieck--Pandharipande \cite{OPand} gives
\[
\sum_n\mathsf P^X_{n,(\beta_h,d)}(-P)^n
=
\sum_g\mathsf N^X_{g,h,d}u^{2g-2},
\qquad P=e^{iu}.
\]
Multiplying by \(Q^{h-1}T^{d-1}\) and summing over \(h,d\) gives
\[
Z^X_{\mathrm{DT}}(T,Q,P)
=\mathcal Z_{\mathrm{OP}}(u,Q,T).
\]
Oberdieck--Pixton \cite[Theorem 1]{OB} proves
\[
\mathcal Z_{\mathrm{OP}}(u,Q,T)
=-\left(\chi_{10}^{\mathrm{OP}}(P,Q,T)\right)^{-1}.
\]
The normalization proposition gives
\[
\chi_{10}^{\mathrm{OP}}(P,Q,T)
=D_5(Z)^2
=4096^{-1}\Delta_5(Z)^2.
\]
Therefore \eqref{eq:dt-op-square} follows.  The sign is fixed by the
common variable \(y=-e^{iu}=-P\) in the pairs comparison and by Bryan's
DT convention \((-p)^n\).
\end{proof}

\subsection{Protected-index replacement}

Replace \texttt{proj.tex:657--669} by:

\begin{conjecture}[Protected D-brane index dictionary]
For the D6--D2--D0 charge
\[
\gamma(h,d,n)=(1;\,(\beta_h,d);\,n),
\]
the physical protected index is expected to be represented by the
Behrend-weighted Hilbert-scheme invariant:
\[
\Omega_X(h,d,n)
=\operatorname{Tr}_{\mathcal H_{\gamma(h,d,n)}}(-1)^F
\stackrel{\mathrm{phys}}{=}
\mathbf{DT}^X_{n,(\beta_h,d)}.
\]
This conjectural dictionary is separate from
Theorem~\ref{thm:op-square}.  The theorem identifies the
quotient-by-\(E\), Behrend-weighted Hilbert-scheme generating series with
the primitive OP curve-counting series.  It does not construct
\(\mathcal H_{\gamma(h,d,n)}\), and it does not identify a microscopic
operator product with the Borcherds bracket.
\end{conjecture}

\subsection{Scalar-branch replacement}

Replace \texttt{proj.tex:672--685} by:

\begin{definition}[Scalar-square branch]
The OP/DT scalar-square branch is the Laurent expansion
\[
Z^X_{\square}
:=Z^X_{\mathrm{OP}}(P,Q,T)
=Z^X_{\mathrm{DT}}(T,Q,P)
=\frac{-4096}{\Delta_5(Z)^2}.
\]
Thus
\[
C_{\square}=C_{\mathrm{OP}}=C_{\mathrm{DT}}=-4096.
\]
If a later argument uses a scalar-square branch from a source other than
the OP/DT theorem, that source must be named separately and may be
written
\[
Z^X_{\square,\mathfrak s}=C_{\square,\mathfrak s}\Delta_5^{-2}.
\]
The symbol \(C_{\mathrm{DT}}\) is not available for such a branch.
\end{definition}

\subsection{Square-root consequence replacement}

Replace \texttt{proj.tex:688--709} by:

\begin{theorem}[Square-root consequence]
For the OP/DT branch,
\[
Z^X_{\square}
=Z^X_{\mathrm{OP}}
=Z^X_{\mathrm{DT}}
=-4096\,\Delta_5(Z)^{-2}.
\]
The normalized Borcherds determinant is
\[
\mathcal D_X(Z)=\Delta_5(Z).
\]
Hence
\[
\mathcal D_X(Z)^2
=\frac{-4096}{Z^X_{\square}},
\qquad
\mathcal Z_{\mathrm{BPS}}^{\mathrm{vir}}:=\mathcal D_X^{-1},
\]
and
\begin{equation}
\label{eq:cell14-square-root}
\left(\mathcal Z_{\mathrm{BPS}}^{\mathrm{vir}}\right)^2
=-\frac{Z^X_{\square}}{4096}
=-\frac{Z^X_{\mathrm{DT}}}{4096}
=-\frac{Z^X_{\mathrm{OP}}}{4096}.
\end{equation}
Equivalently,
\[
\mathcal Z_{\mathrm{BPS}}^{\mathrm{vir}}
=\frac{1}{64}q^{-1/2}r^{-1/2}s^{-1/2}
\sch\mathcal F(\mathcal V).
\]
\end{theorem}

\begin{proof}
The OP/DT theorem gives \(Z^X_{\square}=-4096\Delta_5^{-2}\), and the
normalization theorem gives \(\mathcal D_X=\Delta_5\).  Therefore
\[
(\mathcal D_X^{-1})^2=\Delta_5^{-2}
=-\frac{Z^X_{\square}}{4096}.
\]
The last displayed product is the already normalized inverse
Borcherds determinant.
\end{proof}

\subsection{Status-table replacement}

Replace the four rows in \texttt{proj.tex:720--736} by:

\begin{center}
\small
\begin{tabular}{p{0.25\textwidth}p{0.34\textwidth}p{0.31\textwidth}}
\hline
Object & Equation & Status \\
\hline
Primitive reduced curve counting &
\(Z^X_{\mathrm{OP}}=-4096\Delta_5^{-2}\) &
proved by Oberdieck--Pixton after the normalization
\(\chi_{10}^{\mathrm{OP}}=4096^{-1}\Delta_5^2\) \\
Reduced Hilbert-scheme DT series &
\(Z^X_{\mathrm{DT}}(T,Q,P)=Z^X_{\mathrm{OP}}(P,Q,T)
=-4096\Delta_5^{-2}\) &
proved by Oberdieck's reduced DT/PT theorem, the primitive
Oberdieck--Pandharipande GW/PT comparison, and Oberdieck--Pixton \\
Protected D-brane index &
\(\Omega_X(h,d,n)\stackrel{\mathrm{phys}}{=}
\mathbf{DT}^X_{n,(\beta_h,d)}\) &
conjectural physics dictionary; no microscopic Hilbert space is
constructed here \\
Scalar-square branch &
\(Z^X_{\square}=Z^X_{\mathrm{DT}}=Z^X_{\mathrm{OP}}\),
\(C_{\square}=-4096\) &
theorem for the OP/DT branch; any other scalar-square branch requires
its own named source and constant \\
\hline
\end{tabular}
\end{center}

\subsection{\(C_{\mathrm{DT}}\) replacement ledger}

\begin{center}
\small
\begin{tabular}{p{0.22\textwidth}p{0.68\textwidth}}
\hline
Anchor & Required integration action \\
\hline
\texttt{proj.tex:646--654} &
Delete the Hilbert-scheme DT scalar-normalization assumption and replace
it by Theorem~\ref{thm:op-square} above.  The value becomes
\[
C_{\mathrm{DT}}=-4096.
\]
\\
\texttt{proj.tex:669} &
Remove the reference to \texttt{ass:dt-square}.  The physical
dictionary remains conjectural, but its separation is now from the
enumerative theorem, not from an assumption. \\
\texttt{proj.tex:681--684} &
Replace \(C_{\square,\mathrm{DT}}:=C_{\mathrm{DT}}\) by
\[
C_{\square,\mathrm{DT}}:=-4096.
\]
Remove ``exists under Assumption~\texttt{ass:dt-square}''. \\
\texttt{proj.tex:693--694} &
Replace the two branch lines by
\[
Z^X_{\square}=Z^X_{\mathrm{OP}}=Z^X_{\mathrm{DT}}
\quad\hbox{by Theorem~\ref{thm:op-square}}.
\]
\\
\texttt{proj.tex:724--728} &
Replace the DT table row by the theorem-level row in the
status-table replacement above. \\
\texttt{proj.tex:733--736} &
Replace ``conditional for \(Z^X_{\mathrm{DT}}\)'' by
``theorem for the OP/DT branch''. \\
\texttt{proj.tex:1857--1868,1887--1889} &
Keep \(C_{\square}\) only if the theorem is intentionally about an
arbitrary scalar-square branch.  Add the specialization
\[
\hbox{for the OP/DT branch}\qquad
C_{\square}=-4096,\qquad
(\mathcal Z_{\mathrm{BPS}}^{\mathrm{vir}})^2
=-Z^X_{\mathrm{DT}}/4096.
\]
Do not write \(C_{\mathrm{DT}}\in\mathbb C^\times\). \\
\hline
\end{tabular}
\end{center}

\subsection{Physical separation ledger}

The following manuscript passages must remain physically separated from
the enumerative theorem:
\begin{center}
\small
\begin{tabular}{p{0.24\textwidth}p{0.66\textwidth}}
\texttt{proj.tex:553--556} &
Keep: no microscopic Hilbert space is constructed.\\
\texttt{proj.tex:657--669} &
Replace as above: the protected-index dictionary remains conjectural.\\
\texttt{proj.tex:1278--1280} &
Keep: the denominator algebra is not a microscopic BPS Hilbert space or
operator product.\\
\texttt{proj.tex:1689--1701} &
Keep: index equality is not a state-space identification.\\
\texttt{proj.tex:3397--3418} &
Keep: boundary rows do not inherit CY/DT partition identities or physical
BPS factorization.
\end{tabular}
\end{center}
The integration rule is:
\[
\hbox{DT/PT/OP equality is enumerative;}
\qquad
\hbox{BPS Hilbert-space interpretation is physical and separate.}
\]

\subsection{Commands run}

\begin{itemize}
\item \texttt{nl -ba CLAUDE.md}
\item \texttt{nl -ba /Users/raeez/ecosystem/INVARIANTS.md}
\item \texttt{nl -ba /Users/raeez/ecosystem/AGENTS-HARNESS.md}
\item \texttt{nl -ba proj.tex} around the normalization, scalar-square,
denominator, and physical-separation anchors.
\item \texttt{nl -ba proj.bib} around the OP/DT bibliography block.
\item \texttt{nl -ba agent\_material/10\_reduced\_dt\_pt\_op\_comparison\_theorem.tex}
\item \texttt{nl -ba} on local primary-source TeX files under
\texttt{/tmp/igusa-cell8-sources}.
\item \texttt{rg -n "thm:op-square|ass:dt-square|Protected D-brane index dictionary|Scalar-square branch|Square-root consequence|C\_" proj.tex}
\item Official source lookup: \texttt{https://arxiv.org/abs/1605.04631}
and \texttt{https://link.springer.com/article/10.1007/s00209-017-1953-5}.
\item \texttt{pdflatex -halt-on-error -interaction=nonstopmode} on an
isolated wrapper inputting this dossier.
\item \texttt{rg -n} forbidden-language scan on this dossier.
\item \texttt{git status --short agent\_material/11\_op\_dt\_theorem\_integration\_patch.tex}
\end{itemize}

\subsection{Remaining obligations}

The theorem-level comparison is proved from the cited sources in the
owned dossier.  The remaining obligations are integration obligations:
\begin{enumerate}
\item Insert the bibliography entry in \texttt{proj.bib}.
\item Apply the replacement blocks to \texttt{proj.tex}.
\item Run the manuscript build after integration.
\item Keep generic \(C_{\square}\) only for explicitly named non-OP/DT
branches.
\end{enumerate}
